\documentclass[12pt, a4paper]{article}

\usepackage{fontspec} %Required to change the font.
\usepackage[spanish, es-tabla]{babel}

\usepackage{graphicx} %Required to insert images.
\usepackage{listings} %Required to insert code.
\usepackage[table]{xcolor} %Required to have coloured cells on a table.
\usepackage{amssymb} %Para más símbolos de matemáticas.
\usepackage{amsmath} %Para más símbolos de matemáticas.
\usepackage{amsthm} %Para entornos de teoremas, demostraciones y demás.

% ========== Configuración entornos de teoremas y demás ==========
\newtheorem{teorema}{Teorema}[section]
\newtheorem{lema}[teorema]{Lema}
\newtheorem{proposicion}[teorema]{Proposición}
\newtheorem{corolario}[teorema]{Corolario}

\newtheorem{propiedad}[teorema]{Propiedad}
\newtheorem{axioma}[teorema]{Axioma}
\newtheorem{criterio}[teorema]{Criterio}

\newtheorem{definicion}{Definición}[section]

\newtheorem{ejemplo}{Ejemplo}[section]
\newtheorem{observacion}{Observación}[section]
\newtheorem{nota}{Nota}[section]
\newtheorem{conjetura}{Conjetura}[section]
\newtheorem{problema}{Problema}[section]

% FIN========== Configuración entornos de teoremas y demás ==========

% ========== Entornos para algoritmos en pseudocódigo ==========
\usepackage{algorithm} % Para el entorno "algorithm"
\usepackage{algpseudocode} % Para la sintaxis del pseudocódigo

% Personalización para castellano
\floatname{algorithm}{Algoritmo} % Cambia "Algorithm" por "Algoritmo"
\renewcommand{\algorithmicrequire}{\textbf{Requerir:}}
\renewcommand{\algorithmicensure}{\textbf{Asegurar:}}
\renewcommand{\algorithmicwhile}{\textbf{mientras}}
\renewcommand{\algorithmicif}{\textbf{si}}
\renewcommand{\algorithmicthen}{\textbf{entonces}}
\renewcommand{\algorithmicelse}{\textbf{si no}}
\renewcommand{\algorithmicfor}{\textbf{para}}
\renewcommand{\algorithmicdo}{\textbf{hacer}}
\renewcommand{\algorithmicend}{\textbf{fin}}
\renewcommand{\algorithmicreturn}{\textbf{devolver}}
% FIN========== Entornos para algoritmos en pseudocódigo ==========

\usepackage{geometry} %Para ajustar márgenes
\usepackage{csquotes} %Para citas en español

% ========== Líderes de puntos en el índice ==========
\usepackage{tocloft}

%\renewcommand{\cftchapleader}{\cftdotfill{\cftdotsep}}
\renewcommand{\cftsecleader}{\cftdotfill{\cftdotsep}}
\renewcommand{\cftsubsecleader}{\cftdotfill{\cftdotsep}}
% FIN ========== Líderes de puntos en el índice ==========

% ========== Bibliografía ==========
\usepackage[style=authoryear, backend=biber, language=spanish]{biblatex}
\addbibresource{back_matter/bibliografia.bib}
% FIN ========== Bibliografía ==========


\usepackage{anyfontsize} %Para tamaño de fuente personalizado

% Configuración de márgenes
\geometry{
    left=2.5cm,
    right=2.5cm,
    top=2.5cm,
    bottom=2.5cm
}

%"fontspec" package settings
\setmainfont{Latin Modern Roman}
\setmonofont{Fira Mono}

%"listings" package settings
\lstset{
  language=Swift,
  breaklines=true,
  breakatwhitespace=true,
  postbreak=\mbox{\textcolor{red}{$\hookrightarrow$}\space},
  numbers=left,
  numberstyle=\tiny\ttfamily,
  stepnumber=1,
  numbersep=8pt,
  frame=single,
  lineskip=2pt,
  showstringspaces=false,
  basicstyle=\footnotesize\ttfamily,
  keywordstyle=\color{violet}\bfseries,
  keywordstyle=[2]\color{violet}\bfseries,
  keywordstyle=[3]\color{violet}\bfseries,
  commentstyle=\scriptsize\color{gray}\itshape,
  stringstyle=\color{red},
  identifierstyle=\color{teal},
  directivestyle=\color{green},
  backgroundcolor=\color{gray!5},
  rulecolor=\color{black},
  tabsize=4,
}
\lstdefinestyle{python}{
  language=Python,
  breaklines=true,
  breakatwhitespace=true,
  postbreak=\mbox{\textcolor{red}{$\hookrightarrow$}\space},
  numbers=left,
  numberstyle=\tiny\ttfamily,
  stepnumber=1,
  numbersep=8pt,
  frame=single,
  lineskip=2pt,
  showstringspaces=false,
  basicstyle=\footnotesize\ttfamily,
  keywordstyle=\color{violet}\bfseries,
  keywordstyle=[2]\color{violet}\bfseries,
  keywordstyle=[3]\color{violet}\bfseries,
  commentstyle=\scriptsize\color{gray}\itshape,
  stringstyle=\color{red},
  identifierstyle=\color{teal},
  directivestyle=\color{green},
  backgroundcolor=\color{gray!5},
  rulecolor=\color{black},
  tabsize=4,
}
\renewcommand{\lstlistingname}{Fragmento de código}

% ========== Variables globales ==========
\title{Plantilla de documento corto}
\author{Alfonso Francisco Gil González}
\newcommand{\universidad}{Universidad de Cádiz}
\newcommand{\centro}{Facultad de Ciencias - Escuela Superior de Ingeniería}

\makeatletter
\newcommand{\titulo}{\@title}
\newcommand{\autor}{\@author}
\makeatother

\newcommand{\subtitulo}{Plantilla de \LaTeX para documentos académicos universitarios cortos}
\newcommand{\tipodoc}{Informe de prácticas}
\newcommand{\asignatura}{Sistemas Digitales}
\newcommand{\tutora}{Profesor: Diego Sales Lérida}
\newcommand{\lugar}{Puerto Real}
\newcommand{\fecha}{20 de enero de 2026}
\date{Enero de 2026}
% FIN ========== Variables globales ==========

% ========== Redefinir los tamaños de fuente estándar ==========
\renewcommand{\tiny}{\fontsize{8pt}{9.6pt}\selectfont} % 8pt, interlineado 9.6pt
\renewcommand{\scriptsize}{\fontsize{12pt}{14.4pt}\selectfont}
\renewcommand{\footnotesize}{\fontsize{10pt}{12pt}\selectfont}
\renewcommand{\small}{\fontsize{10pt}{12pt}\selectfont}
\renewcommand{\normalsize}{\fontsize{12pt}{14.4pt}\selectfont}
\renewcommand{\large}{\fontsize{14pt}{16.8pt}\selectfont}
\renewcommand{\Large}{\fontsize{16pt}{19.2pt}\selectfont}
\renewcommand{\LARGE}{\fontsize{20pt}{24pt}\selectfont}
\renewcommand{\huge}{\fontsize{22pt}{26.4pt}\selectfont}
\renewcommand{\Huge}{\fontsize{24pt}{28.8pt}\selectfont}
% FIN ========== Redefinir los tamaños de fuente estándar ==========

% ========== Encabezados y pies de páginas ==========
\usepackage{fancyhdr}
\pagestyle{fancy}

% Limpia los estilos por defecto
\fancyhf{}

% Configura el encabezado
\fancyhead[L]{\small \autor}
\fancyhead[R]{\small \titulo}

% Configura el pie de página
\fancyfoot[C]{\small \thepage}

\renewcommand{\headrulewidth}{0pt} % Elimina la línea del encabezado

% Pies de página del page style plain, usado al iniciar un capítulo
\fancypagestyle{plain}{
    \fancyhf{}
    \fancyfoot[C]{\small \thepage}
}
% FIN ========== Encabezados y pies de páginas ==========

\usepackage[hidelinks]{hyperref} %Required to make "Contents" clickable.
\usepackage{bookmark} %Improves "hyperref" package.
\usepackage{cleveref}

% ========== Configuración "cleveref" ==========
\crefname{teorema}{teorema}{teoremas}
\Crefname{teorema}{Teorema}{Teoremas}
\crefname{lema}{lema}{lemas}
\Crefname{lema}{Lema}{Lemas}
\crefname{proposicion}{proposición}{proposiciones}
\Crefname{proposicion}{Proposición}{Proposiciones}
\crefname{corolario}{corolario}{corolarios}
\Crefname{corolario}{Corolario}{Corolarios}
\crefname{propiedad}{propiedad}{propiedades}
\Crefname{propiedad}{Propiedad}{Propiedades}
\crefname{axioma}{axioma}{axiomas}
\Crefname{axioma}{Axioma}{Axiomas}
\crefname{criterio}{criterio}{criterios}
\Crefname{criterio}{Criterio}{Criterios}
\crefname{definicion}{definición}{definiciones}
\Crefname{definicion}{Definición}{Definiciones}
\crefname{ejemplo}{ejemplo}{ejemplos}
\Crefname{ejemplo}{Ejemplo}{Ejemplos}
\crefname{observacion}{observación}{observaciones}
\Crefname{observacion}{Observación}{Observaciones}
\crefname{nota}{nota}{notas}
\Crefname{nota}{Nota}{Notas}
\crefname{conjetura}{conjetura}{conjeturas}
\Crefname{conjetura}{Conjetura}{Conjeturas}
\crefname{problema}{problema}{problemas}
\Crefname{problema}{Problema}{Problemas}
\crefname{algorithm}{algoritmo}{algoritmos}
\Crefname{algorithm}{Algoritmo}{Algoritmos}
\crefname{listing}{fragmento de código}{fragmentos de código}
\Crefname{listing}{Fragmento de código}{Fragmentos de código}
\crefname{equation}{ecuación}{ecuaciones}
\Crefname{equation}{Ecuación}{Ecuaciones}
% FIN========== Configuración "cleveref" ==========


\begin{document}

% Front matter
\pagenumbering{roman}
\setcounter{page}{0}

\thispagestyle{empty}
\null %Para que el centrado vertical con los "vfill" funcione correctamente
\vfill
\begin{center}
{\huge \universidad \par}
\vspace{0.5cm}
{\LARGE \centro \par}
\vspace{2.5cm}
{\textbf{\Huge \titulo} \par}
\vspace{1cm}
{\LARGE \subtitulo \par}
\vspace{2.5cm}
{\Large \tipodoc \par}
\vspace{0.5cm}
{\Large \asignatura \par}
\vspace{2.5cm}
{\LARGE \autor \par}
\vspace{1cm}
{\LARGE \tutora \par}
\vspace{2.5cm}
{\large \lugar \par}
\vspace{0.5cm}
{\large \fecha \par}
\end{center}
\vfill

\chapter*{Resumen}
\addcontentsline{toc}{chapter}{Resumen}
Lorem ipsum
\clearpage

\tableofcontents


% Main matter
\clearpage
\pagenumbering{arabic}
\setcounter{page}{1}

\chapter*{Introducción}
\addcontentsline{toc}{chapter}{Introducción}
Lorem ipsum 

Lorem ipsum 

Lorem ipsum 

Lorem ipsum 

Lorem ipsum 

Lorem ipsum 

Lorem ipsum 

Lorem ipsum 

Lorem ipsum 

Lorem ipsum 

Lorem ipsum 

Lorem ipsum 

Lorem ipsum 

Lorem ipsum 

Lorem ipsum 

Lorem ipsum 

Lorem ipsum 

Lorem ipsum 

Lorem ipsum 

Lorem ipsum 

Lorem ipsum 

Lorem ipsum 

Lorem ipsum 

Lorem ipsum 

Lorem ipsum 

Lorem ipsum 

Lorem ipsum 

Lorem ipsum 

Lorem ipsum 

Lorem ipsum 

Lorem ipsum 

Lorem ipsum 

Lorem ipsum 

Lorem ipsum 

Lorem ipsum 

Lorem ipsum 

Lorem ipsum 

Lorem ipsum 

Lorem ipsum 

Lorem ipsum 

Lorem ipsum 

Lorem ipsum 

Lorem ipsum 

Lorem ipsum 

Lorem ipsum 

Lorem ipsum 

Lorem ipsum 

Lorem ipsum 

Lorem ipsum 

Según \textcite{apellido2023}, el método es eficiente.


\clearpage
\chapter{Título del primer capítulo}
Lorem ipsum

\[
  \int_0^{\mathrm{\pi}} \sin x \, \mathrm{d}x = 2
\]

\clearpage
\chapter{Título del segundo capítulo}
Lorem ipsum

\clearpage
\section*{Conclusión}
\addcontentsline{toc}{section}{Conclusión}

Lorem ipsum dolor sit amet, consectetur adipiscing elit, sed do eiusmod tempor incididunt ut labore et dolore magna aliqua. Ut enim ad minim veniam, quis nostrud exercitation ullamco laboris nisi ut aliquip ex ea commodo consequat. Duis aute irure dolor in reprehenderit in voluptate velit esse cillum dolore eu fugiat nulla pariatur. Excepteur sint occaecat cupidatat non proident, sunt in culpa qui officia deserunt mollit anim id est laborum.

Lorem ipsum.

\begin{teorema}
    Contenido del teorema.
\end{teorema}

\begin{proof}
    Contenido de la demostración.
\end{proof}

\begin{corolario}
    Contenido del corolario.
\end{corolario}

\begin{proof}
    Contenido de la demostración.
\end{proof}

\begin{proposicion}
    Contenido de la proposición.
\end{proposicion}

\subsection*{Título de la subsección}
\addcontentsline{toc}{subsection}{Título de la subsección}

Lorem ipsum dolor sit amet, consectetur adipiscing elit, sed do eiusmod tempor incididunt ut labore et dolore magna aliqua. Ut enim ad minim veniam, quis nostrud exercitation ullamco laboris nisi ut aliquip ex ea commodo consequat. Duis aute irure dolor in reprehenderit in voluptate velit esse cillum dolore eu fugiat nulla pariatur. Excepteur sint occaecat cupidatat non proident, sunt in culpa qui officia deserunt mollit anim id est laborum.

Lorem ipsum.

\begin{teorema}
    Contenido del teorema.
\end{teorema}

\begin{proof}
    Contenido de la demostración.
\end{proof}

\begin{corolario}
    Contenido del corolario.
\end{corolario}

\begin{proof}
    Contenido de la demostración.
\end{proof}

\begin{proposicion}
    Contenido de la proposición.
\end{proposicion}

\subsubsection*{Título de la subsubsección}
\addcontentsline{toc}{subsubsection}{Título de la subsubsección}

Lorem ipsum dolor sit amet, consectetur adipiscing elit, sed do eiusmod tempor incididunt ut labore et dolore magna aliqua. Ut enim ad minim veniam, quis nostrud exercitation ullamco laboris nisi ut aliquip ex ea commodo consequat. Duis aute irure dolor in reprehenderit in voluptate velit esse cillum dolore eu fugiat nulla pariatur. Excepteur sint occaecat cupidatat non proident, sunt in culpa qui officia deserunt mollit anim id est laborum.

Lorem ipsum.

\begin{teorema}
    Contenido del teorema.
\end{teorema}

\begin{proof}
    Contenido de la demostración.
\end{proof}

\begin{corolario}
    Contenido del corolario.
\end{corolario}

\begin{proof}
    Contenido de la demostración.
\end{proof}

\begin{proposicion}
    Contenido de la proposición.
\end{proposicion}


% Back matter
\clearpage
\printbibliography[title=Bibliografía]
\addcontentsline{toc}{section}{Bibliografía}

\clearpage
\section*{Anexo 1}
\addcontentsline{toc}{section}{Anexo 1}
Lorem ipsum

\clearpage
\chapter*{Anexo 2}
\addcontentsline{toc}{chapter}{Anexo 2}
Lorem ipsum


\end{document}